\section{Digital Audio Restoration: Restoration of Local Degradations}

\subsection{Introduction}
Historical recordings are often affected by defects of finite duration, which occur at random locations and only affect the signal in a localized manner.
They are perceived in different ways. A common manifestation is the presence of tiny tick noises, which may occur in any recording medium, including modern digital sources.
Another typical case is scratch and crackle noise, usually associated with analog disc recordings.

For example, a poor-quality 78 rpm record might typically exhibit around 2000 clicks per second, with durations ranging from less than 20 $\mu$s up to 4 ms in extreme cases.
Nevertheless, in most cases at least 90\% of the samples remain undegraded, so it is reasonable to hope that a convincing restoration can be achieved.

The causes of such degradations are diverse. They may result from specks of dirt or dust adhering to the grooves of a disc, from the granularity of the material used for pressing, or from mechanical damage such as surface scratches.

\begin{figure}[H]
    \centering
    \includegraphics[width=0.35\linewidth]{Phonographs degradation causes.png}
    \caption{Electron micrograph showing dust and damage to the grooves of a 78 rpm grammophone disc}
\end{figure}

Let us take as example a \textbf{click-degraded music waveform}.
In the time domain, we find sharp impulsive peaks on the waveform, caused by the clicks.
In the frequency domain, the sonogram reveals how these clicks manifest as vertical lines across the time-frequency representation.

\begin{figure}[H]
    \centering
    \includegraphics[width=0.8\linewidth]{Click degredation.png}
    \caption{Click-degraded music waveform from a 78 rpm recording: time-domain waveform (left) and time-frequency sonogram (right)}
\end{figure}

\subsection{Modelling clicks}
A defect in an audio signal may manifest itself either by adding unwanted noise to the signal, or by replacing the original waveform for a short duration of time, defining two conceptual models.
The \textbf{additive model} is generally acceptable for most surface defects in recording media, including small scratches, dust or dirt.
The \textbf{replacement model} is more appropriate for severe degradations such as very large scratches or breakages, which completely obliterate the underlying signal information.
Such cases often excite long-term resonances in mechanical playback systems and therefore need to be treated differently.

Let us consider the additive model, having the underlying clean audio $s$, the corrupting noise process $v$ and an indicator function $i$:

\[
x(n) = s(n) + i(n)\ v(n),
\qquad
i(n) =
\begin{cases}
    1, & \text{when the degradation is present} \\
    0, & \text{otherwise}
\end{cases}
\]

Click degradations often appear as \textbf{contiguous bursts of corrupted samples}, which start at random positions and last for a random duration, typically between 1 and 200 samples at a sampling rate of 44.1 kHz.

There is a \textbf{strong dependence} between successive samples of the \textbf{switching process} $i$, which implies that the noise cannot simply be assumed as impulsive.
Treating \textbf{clusters of impulses} is significantly more difficult than dealing with isolated impulses.
Furthermore, adjacent impulses may cancel each other in the detection space (leading to missed detections), or may add constructively (producing the impression of more impulses and thus leading to false detections).  

The \textbf{amplitude} of the clicks is also very \textbf{variable}.
On one hand, the largest click amplitudes can exceed the largest amplitudes of the underlying signal.
On the other hand, the smallest audible defects may lie more than 40 dB below the local signal level.  

Several \textbf{difficulties} arise in detection and restoration.
Large-amplitude defects tend to bias parameter estimations and threshold determination procedures, which may result in smaller defects remaining undetected.
Moreover, the choice of an \textbf{appropriate threshold} for detection schemes becomes a delicate and difficult problem in practice.

\subsection{Approach to click restoration}
In click restoration, two main tasks must be carried out.
The first task is \textbf{detection}, which consists in estimating the switching process $i$.
This step allows us to decide which samples of the audio signal are degraded and thus affected by impulsive noise.

The second task is \textbf{signal estimation}, which aims to reconstruct the underlying audio data whenever corruption is present.
In the absence of direct information about the true, uncorrupted signal, this reconstruction is typically achieved by exploiting the information contained in neighbouring uncorrupted samples.

\begin{figure}[H]
    \centering
    \includegraphics[width=0.5\linewidth]{Click restoration scheme.png}
    \caption{Click restoration approach}
\end{figure}

The simplest approach to \textbf{detect clicks} is to magnify them with respect to the signal by applying a \textbf{high-pass filter}, then performing thresholding.
This strategy is straightforward to implement and does not require any knowledge of the system parameters.

However, the detection relies on the \textbf{assumption} that there is a clear \textbf{frequency separation} between the signal and the clicks, which is rarely the case in practice.
Consequently, this method suffers when the signal itself has a strong high-frequency content, or when the clicks are band-limited.  
This highlights the strong variability of click characteristics and the limitations of such a basic detection approach.

\begin{figure}[H]
    \centering
    \includegraphics[width=0.5\linewidth]{Click detection.png}
    \caption{Simplest click detection approach}
\end{figure}

One possible approach to \textbf{estimate the click-degraded signal} is to use a \textbf{median filter}: the corrupted samples are replaced by values that are statistically likely, obtained from the median of their neighbouring samples.  
This approach typically causes only a modest degradation of surrounding, uncorrupted samples.  
However, it becomes unsuitable when the \textbf{gaps} caused by the degradation are larger than just a \textbf{few samples}, since the median filter is unable to reconstruct longer stretches of missing data.

Another possible approach is to \textbf{splice uncorrupted data} from nearby regions into the gap.  
This method can be effective but requires careful handling of discontinuities that may arise at the boundaries between the spliced and original signal.  
In practice, splicing often makes use of techniques borrowed from time-scaling and pitch-scaling algorithms, such as SOLA or PSOLA.
For this approach to work convincingly, it is important that the underlying signal is \textbf{periodic} (or at least quasi-periodic) and that the pitch period is estimated with sufficient accuracy.

\subsubsection{Median filter for signal estimation}
A median filter operates by scanning through the signal entry by entry, and replacing each value with the \textbf{median} of its neighbouring entries.  
The median is defined as the middle value of the sorted entries within the chosen window.  
To make this operation straightforward, the filter is usually applied with a sliding window of odd length, since this guarantees a well-defined central value.

For example, let us consider the following signal:

\[
x = \{2, 80, 6, 3, 5, 4, \cdots\}
\]

We want to apply a median filter of length 3.  
At the boundaries, the window is completed by repeating the nearest available entries.  
The operation to obtain the median-filtered signal $y$ proceeds as follows:

\[
\left.
\begin{array}{ccc}
y(1) =& \mathrm{Median}\{2,2,80\} &= 2 \\
y(2) =& \mathrm{Median}\{2,6,80\} &= 6 \\
y(3) =& \mathrm{Median}\{80,6,3\} &= 6 \\
y(4) =& \mathrm{Median}\{6,3,5\} &= 5 \\
y(5) =& \mathrm{Median}\{3,5,4\} &= 4
\end{array} \qquad
\right\} \quad \Rightarrow \quad
y = \{2, 6, 6, 5, 4, \cdots\}
\]

This simple example illustrates how the median filter suppresses outliers (such as the value $80$) while preserving the general structure of the surrounding signal.

Strictly speaking, the median filter is \textbf{not a true filter} in the usual sense, because it is not a linear time-invariant (LTI) system.
Its non-linearity arises from the \textbf{sorting operation} used to determine the median.

\begin{figure}[H]
    \centering
    \includegraphics[width=0.3\linewidth]{Median filter.png}
    \caption{Median filter}
\end{figure}

One of the main \textbf{advantages} of the median filter is that it \textbf{preserves step-like functions}.
This property makes it particularly useful when dealing with signals or images containing sharp transitions (such as edges) where conventional linear filters would blur the details.
Another important feature is its ability to \textbf{remove isolated errors} or outliers.
For example, a filter of length $2N+1$ can eliminate anomalous sequences of up to $N$ samples by replacing them with values consistent with their neighbours.

\begin{figure}[H]
    \centering
    \includegraphics[width=0.6\linewidth]{Median vs linear filter.png}
    \caption{Comparison between linear and 5-median filters for different input signals}
\end{figure}

However, the method has \textbf{limitations}.
Median filtering is suitable \textbf{only for very short bursts} of missing or corrupted data, typically spanning just a few samples.
If the gaps become longer, the signal degradation becomes unacceptable and fine local details are quickly lost.
This trade-off means that while the filter is excellent for impulsive noise removal, it is less effective when larger sections of data are corrupted.

\subsection{General scheme for detection and estimation}
\textbf{Blind approaches} to click detection and signal estimation are usually not very effective.  
Instead, we need to rely on \textbf{model-based methods}, such as sinusoidal modelling, autoregressive (AR) modelling or periodic models.  

In principle, the optimal solution would be to carry out joint detection and estimation.  
However, this is a very challenging problem, because the two steps are not independent of each other.  
One way to address this difficulty is through an iterative refinement approach.

We decide to adopt an \textbf{autoregressive (AR) model} to represent the signal:

\[
s(n) = \sum_{k=1}^P a_k \ s(n-k) + e(n)
\]

The parameters $a_k$ are estimated from the available signal, while $e(n)$ is a residual error term.

The iterative procedure works as follows: given the current indicator function $i(n)$, which specifies which samples are corrupted, we first estimate the AR parameters $a_k$.  
These parameters are then used to improve the detection of $i(n)$.  
If the newly obtained indicator function differs from the one of the previous iteration, the process is repeated.  
Once the detection stabilizes, signal estimation is performed: using both $i(n)$ and $a_k$, the missing or corrupted samples are reconstructed from the model, thereby filling the gaps.

\begin{figure}[H]
    \centering
    \includegraphics[width=0.8\linewidth]{Detection estimation scheme.png}
    \caption{General scheme for detection and estimation in click degradation problems}
\end{figure}

This scheme provides a structured way to alternate between parameter estimation, detection, and signal reconstruction, converging toward a more reliable restoration of the audio.

\subsubsection{Click detection}
We consider the audio data $s$ to be drawn from a short-term stationary AR process:

\[
s(n) = \sum_{k=1}^{P} a_k \ s(n-k) + e(n)
\]

For now, we assume that the AR model parameters $a_k$ and the excitation variance $\sigma_e^2$ are known.

\begin{figure}[H]
    \centering
    \includegraphics[width=0.35\linewidth]{Click detection block.png}
    \caption{Click detection block}
\end{figure}

Applying the $\mathcal{Z}$-transform, the AR model becomes:

\[
s(n) = \sum_{k=1}^{P} a_k \ s(n-k) + e(n) 
\quad \xrightarrow{\ \mathcal{Z}\ }\quad 
S(z) = S(z)\sum_{k=1}^{P} a_k \ z^{-k} + E(z)
\]

From this, the \textbf{prediction error filter} $H$ is defined as:

\[
H(z) = \frac{E(z)}{S(z)} = 1 - \sum_{k=1}^{P} a_k \ z^{-k}
\]

If we feed the prediction error filter with the corrupted data $x$:

\[
e_d(n) = x(n) - \sum_{k=1}^P a_k \ x(n-k)
= \left[s(n) + i(n)v(n)\right] - \sum_{k=1}^P a_k \left[s(n-k) + i(n-k) v(n-k)\right]
\]

\[
= \underbrace{s(n) - \sum_{k=1}^P a_k \ s(n-k)}_{e(n)} + \underbrace{i(n) v(n) - \sum_{k=1}^P a_k \ i(n-k) v(n-k)}_{u \ast h \ (n)}
\]

We obtain an expression with $h$ the prediction error filter impulse response and $u(n)=i(n)v(n)$ the localized noise.

\begin{figure}[H]
    \centering
    \includegraphics[width=0.7\linewidth]{Click detection 2.png}
    \caption{Prediction error filter}
\end{figure}

The response of the prediction error filter to the corrupted signal is essentially the noise component $e$.  
The filter suppresses the predictable part of the musical signal while leaving the unpredictable component $v$ almost unaffected.  

We remind the definition of the AR process $s$:

\[
s(n) = \sum_{k=1}^P a_k \ s(n-k) + e(n)
\]

It is intuitive to think that the signal $s$ may be obtained from the error $e$ using the \textbf{inverse process} of the prediction filtering $H$.

\begin{figure}[H]
    \centering
    \includegraphics[width=0.35\linewidth]{Click detection 3.png}
    \caption{Inverse of the prediction error filter}
\end{figure}

This means that the variance $\sigma_e$ of the prediction error $e$ can be directly related to the variance $\sigma_s$ of the uncorrupted signal $s$.  
Specifically, the power spectral density satisfies the following:

\[
P_s(e^{j\omega}) = P_e(e^{j\omega}) \abs{\frac{1}{H(e^{j\omega})}}^2
\quad \Rightarrow \quad
2\pi \sigma_s^2 = \int_{-\pi}^{\pi} P_e(e^{j\omega}) \abs{\frac{1}{H(e^{j\omega})}}^2 d\omega
= \sigma_e^2 \int_{-\pi}^{\pi} \abs{\frac{1}{H(e^{j\omega})}}^2 d\omega
\]

The \textbf{variance} of the excitation (i.e. prediction error) is:

\[
\sigma_e^2 = \frac{2\pi \sigma_s^2}{\int_{-\pi}^{\pi} \abs{ \frac{1}{H(e^{j\omega})} }^2 d\omega}
\]

The variance expression we found can be used to quantify the reduction in the power of the audio signal.  
For highly correlated audio, this reduction is significant, often reaching values of 40 dB or more.  

We remind the prediction error $e_d$ expression:

\[
e_d(n) = e(n) + i(n)v(n) - \sum_{k=1}^P a_k \ i(n-k) v(n-k)
\]

This relation tells us that if a click occurs as a single impulse, it triggers the full impulse response of the prediction error filter, scaled by the magnitude of the impulse.  
As a result, localized noise-like disturbances are strongly amplified when they are uncorrelated with the underlying signal.  
This \textbf{amplification}, however, comes at the \textbf{cost of temporal localization}, since the effect of a single click extends across $P$ samples.

Click detection is carried out by applying a \textbf{threshold} to the \textbf{prediction error} $e_d$.  
The threshold is chosen based on the magnitude of the peaks and on the order of the prediction model, which introduces a trade-off between false positives and false negatives.

\begin{figure}[H]
    \centering
    \includegraphics[width=0.45\linewidth]{Click detection thresholding.png}
    \caption{Click detection: original signal (top) and absolute value of prediction error (bottom)}
\end{figure}

The procedure can be summarized by the following algorithm:  

\begin{enumerate}
    \item For $n = 1$ to $N$, compute the prediction error:
    
    \[
    e(n) = x(n) - \sum_{i=1}^{P} a_i \ x(n-i)
    \]

    \item Set $i(n)$ accordingly:

    \[
    i(n) =
    \begin{cases}
        1, & \abs{e(n)} > k \ \sigma_e \\
        0, & \abs{e(n)} \leq k \ \sigma_e
    \end{cases}
    \] 
\end{enumerate}

The parameter $k$ controls the balance between missed detections and false alarms, and its choice directly impacts the accuracy of the detection process.

\subsubsection{Parameter estimation}
We now turn our attention to the problem of modelling \textbf{parameter estimation}.  
The estimation task is simplified by the fact that the model is \textbf{linear in the parameters}, which makes the computations more tractable.  
Moreover, the AR model is assumed to be a valid representation of the signal dynamics.  
Under this assumption, the \textbf{prediction error} (or innovation) behaves as a white process, meaning that the samples are \textbf{uncorrelated and identically distributed}.  
In addition, we consider the prediction error to be \textbf{Gaussian}, which implies that uncorrelatedness is equivalent to independence.  

\begin{figure}[H]
    \centering
    \includegraphics[width=0.35\linewidth]{Parameter estimation block.png}
    \caption{Parameter estimation block}
\end{figure}

Let us consider a random process $x(n)$ generated by an order-$P$ autoregressive (AR) filter driven by a white noise innovation process $e(n)$.
The AR model is the following:

\[
x(n) = \sum_{k=1}^{P} a_k \ x(n-k) + e(n)
\quad \Rightarrow \quad
e(n) = x(n) - \sum_{k=1}^{P} a_k \ x(n-k)
\]

We denote as $p_e(e(n))$ the Probability Density Function (PDF) of the innovation process.
The probability of observing the next sample of the process given the previous $P$ samples is the following:

\[
P\left[x(n) \ \mid \ x(n-1),x(n-2),\cdots,x(n-P)\right] = p_e \left( \underbrace{x(n) - \sum_{k=1}^{P} a_k \ x(n-k)}_{e(n)} \right)
\]

Since the samples of $e(n)$ are assumed to be independent and identically distributed (i.i.d.), the joint probability of a sequence of samples from $x(P+1)$ up to $x(N)$ given the initial $P$ samples is given by:

\[
P\left[x(P+1),x(P+2),\dots,x(N) \ \mid \ x(1),x(2),\dots,x(P)\right] 
= \prod_{n=P+1}^{N} p_e \left( x(n) - \sum_{k=1}^{P} a_k x(n-k) \right)
\]

Let us define the model in matrix form:

\[
\underline{\mathrm{x}} =
\begin{bmatrix}
x(1) \\[3pt]
\vdots \\[3pt]
x(N)
\end{bmatrix}
=
\begin{bmatrix}
\underline{\bar{\mathrm{x}}} \\[3pt]
\underline{\hat{\mathrm{x}}}
\end{bmatrix},
\qquad
\underline{\bar{\mathrm{x}}} =
\begin{bmatrix}
x(1) \\[3pt]
\vdots \\[3pt]
x(P)
\end{bmatrix},
\qquad
\underline{\hat{\mathrm{x}}} =
\begin{bmatrix}
x(P+1) \\[3pt]
\vdots \\[3pt]
x(N)
\end{bmatrix}
\]

We rewrite the AR model as follows:

\[
x(n) = \sum_{k=1}^{P} a_k x(n-k) + e(n)
\quad \Rightarrow \quad
\underline{\hat{\mathrm{x}}} =
\begin{bmatrix}
x(P+1) \\[3pt]
\vdots \\[3pt]
x(N)
\end{bmatrix}
=
\begin{bmatrix}
\sum_{k=1}^{P} a_k \ x(P+1-k) \\[3pt]
\vdots \\[3pt]
\sum_{k=1}^{P} a_k \ x(N-k)
\end{bmatrix}
+
\begin{bmatrix}
e(P+1) \\[3pt]
\vdots \\[3pt]
e(N)
\end{bmatrix}
\]

Equivalently, we can express this as the following \textbf{general linear model}:

\[
\underline{\hat{\mathrm{x}}} =
\begin{bmatrix}
x(P+1) \\[3pt]
\vdots \\[3pt]
x(N)
\end{bmatrix}
=
\underbrace{
\begin{bmatrix}
x(P) & x(P-1) & \cdots & x(1) \\[3pt]
x(P+1) & x(P) & \cdots & x(2) \\[3pt]
\vdots & \vdots & \ddots & \vdots \\[3pt]
x(N-1) & x(N-2) & \cdots & x(N-P)
\end{bmatrix}}_{\left[\mathrm{G_x}\right]}
\underbrace{
\begin{bmatrix}
a_1 \\[3pt]
\vdots \\[3pt]
a_P
\end{bmatrix}}_{\underline{\mathrm{a}}}
+
\underbrace{
\begin{bmatrix}
e(P+1) \\[3pt]
\vdots \\[3pt]
e(N)
\end{bmatrix}}_{\underline{\mathrm{e}}}
\quad \Rightarrow \quad
\underline{\hat{\mathrm{x}}} = \left[\mathrm{G_x}\right] \ \underline{\mathrm{a}} + \underline{\mathrm{e}}
\]


This means that the previous conditional probability:

\[
P \left[ x(P+1), x(P+2), \ldots, x(N) \ \mid \ x(1), x(2), \ldots, x(P) \right]
= \prod_{n=P+1}^{N} p_e \left( x(n) - \sum_{k=1}^{P} a_k \ x(n-k) \right)
\]

Can be rewritten as:

\[
P[\underline{\hat{\mathrm{x}}} \mid \underline{\bar{\mathrm{x}}}, \ \underline{\mathrm{a}}] = p_e \left( \underline{\hat{\mathrm{x}}} - \left[\mathrm{G_x}\right] \ \underline{\mathrm{a}} \right)
\]

In the additional case of Gaussian excitation with zero-mean and variance $\sigma_e^2$:

\[
p_e \left( \underline{\hat{\mathrm{x}}} - \left[\mathrm{G_x}\right] \ \underline{\mathrm{a}} \right)
= \left( \frac{1}{2\pi\sigma_e^2} \right)^{\tfrac{N-P}{2}}
\exp \left(
-\frac{1}{2\sigma_e^2} \ \left( \underline{\hat{\mathrm{x}}} - \left[\mathrm{G_x}\right] \ \underline{\mathrm{a}} \right)^T \left( \underline{\hat{\mathrm{x}}} - \left[\mathrm{G_x}\right] \ \underline{\mathrm{a}} \right)
\right)
\]

Let us recall the problem: we want to estimate the AR model parameters $\underline{\mathrm{a}}$ from the observed data $\underline{\mathrm{x}}$.
The \textbf{Maximum Likelihood (ML) estimator} treats the parameters as unknown constants with no prior information.
The observed data $\underline{\mathrm{x}}$ is considered random and we can often obtain the PDF for $\underline{\mathrm{x}}$ when the value of $\underline{\mathrm{a}}$ is known.
This PDF is termed \textbf{likelihood} $L(\underline{\mathrm{x}}; \underline{\mathrm{a}})$ and is defined as follows:

\[
L(\underline{\mathrm{x}}; \underline{\mathrm{a}}) 
= P\left[\underline{\hat{\mathrm{x}}} \mid \underline{\bar{\mathrm{x}}}, \ \underline{\mathrm{a}}\right] 
= p_e \left(\underline{\hat{\mathrm{x}}} - \left[\mathrm{G_x}\right] \ \underline{\mathrm{a}}\right)
\]

The ML estimate for $\underline{\mathrm{a}}$ is that value $\underline{\mathrm{a}}^\mathrm{ML}$ that maximises the likelihood for given observations $\underline{\mathrm{x}}$:

\[
\underline{\mathrm{a}}^\mathrm{ML} = \arg \max_{\underline{\mathrm{a}}} \left\{ L(\underline{\mathrm{x}}; \underline{\mathrm{a}}) \right\}
\]

The rationale behind this is that the ML solution corresponds to the parameter vector which would have generated the observed data $\underline{\mathrm{x}}$ with highest probability.  
It is usually computed by maximizing the log-likelihood.

In our case, the ML estimator is obtained as:

\[
\underline{\mathrm{a}}^\mathrm{ML} = \arg \max_{\underline{\mathrm{a}}} \left\{ \log \left( p_e \left( \underline{\hat{\mathrm{x}}} - \left[\mathrm{G_x}\right]\underline{\mathrm{a}} \right) \right) \right\}
= \arg \max_{\underline{\mathrm{a}}} \left\{ -\frac{1}{2\sigma_e^2} \ \underline{\mathrm{e}}^T \underline{\mathrm{e}} \right\}
= \arg \min_{\underline{\mathrm{a}}} \left\{ \underline{\mathrm{e}}^T \underline{\mathrm{e}} \right\}
= \arg \min_{\underline{\mathrm{a}}} E,
\qquad
E = \underline{\mathrm{e}}^T \underline{\mathrm{e}}
\]

This result can be found by setting the following condition:

\[
\frac{\partial E}{\partial \underline{\mathrm{a}}}
= 2 \underline{\mathrm{e}}^T \frac{\partial \underline{\mathrm{e}}}{\partial \underline{\mathrm{a}}}
= 2 \left( \underline{\hat{\mathrm{x}}} - \left[\mathrm{G_x}\right]\underline{\mathrm{a}} \right)^T 
\frac{\partial \left(\underline{\hat{\mathrm{x}}} - \left[\mathrm{G_x}\right]\underline{\mathrm{a}}\right)}{\partial \underline{\mathrm{a}}}
= 2 \left( \underline{\hat{\mathrm{x}}} - \left[\mathrm{G_x}\right] \ \underline{\mathrm{a}} \right)^T \left[\mathrm{G_x}\right]
= 0
\]

Finally, this can be solved with respect to the parameter vector $\underline{\mathrm{a}}$ as:

\[
\underline{\mathrm{a}}^\mathrm{ML} =
\left( \left[\mathrm{G_x}\right]^T \left[\mathrm{G_x}\right] \right)^{-1} 
\left[\mathrm{G_x}\right]^T \underline{\hat{\mathrm{x}}}
\]

\subsubsection{Signal estimation}
We now focus on the problem of estimating the missing data: we can highlight some useful facts and hypotheses.  
The model is \textbf{linear with respect to the data}, which simplifies the estimation.  
The AR model is suitable for capturing the dynamics of the clean audio signal.
As before, the prediction error (i.e. innovation) is assumed to be a white process, meaning it is uncorrelated and identically distributed across samples, and the prediction error is Gaussian, so incorrelation implies independency.

\begin{figure}[H]
    \centering
    \includegraphics[width=0.35\linewidth]{Signal estimation block.png}
    \caption{Signal estimation block}
\end{figure}

\textbf{Least-Square Auto-Regressive (LSAR) interpolation} is usually employed in digital click removal, where samples are obliterated.
We start from a block of data samples $x(n)$, assumed to be drawn from an AR process with coefficients $a_k$.

\begin{figure}[H]
    \centering
    \includegraphics[width=0.65\linewidth]{LSAR.png}
    \caption{LSAR interpolation: known samples $\underline{\bar{\mathrm{x}}}$ and unknown samples $\underline{\hat{\mathrm{x}}}$ of the gap}
\end{figure}

To process the data, we work in a \textit{batch mode} by partitioning it into three blocks: $\underline{\bar{\mathrm{x}}}_L$ known samples on the left, $\underline{\hat{\mathrm{x}}}$ unknown samples (the gap), $\underline{\bar{\mathrm{x}}}_R$ known samples on the right.

\[
\underline{\bar{\mathrm{x}}}_L =
\begin{bmatrix}
x(1) \\[3pt]
\vdots \\[3pt]
x(m)
\end{bmatrix},
\qquad
\underline{\hat{\mathrm{x}}} =
\begin{bmatrix}
x(m+1) \\[3pt]
\vdots \\[3pt]
x(m+l)
\end{bmatrix},
\qquad
\underline{\bar{\mathrm{x}}}_R =
\begin{bmatrix}
x(m+l+1) \\[3pt]
\vdots \\[3pt]
x(N)
\end{bmatrix}
\]

The complete data block can be written as follows:

\[
\underline{\mathrm{x}} =
\begin{bmatrix}
\underline{\bar{\mathrm{x}}}_L \\[3pt]
\underline{\hat{\mathrm{x}}} \\[3pt]
\underline{\bar{\mathrm{x}}}_R
\end{bmatrix},
\qquad
\underline{\bar{\mathrm{x}}} =
\begin{bmatrix}
\underline{\bar{\mathrm{x}}}_L \\[3pt]
\underline{\bar{\mathrm{x}}}_R
\end{bmatrix}: \ \text{known samples}
\]

Let us rearrange the matrices as follows:

\[
\underline{\mathrm{x}} =
\begin{bmatrix}
\underline{\bar{\mathrm{x}}}_L \\[3pt]
\underline{\hat{\mathrm{x}}} \\[3pt]
\underline{\bar{\mathrm{x}}}_R
\end{bmatrix}_{N \times 1}
=
\underbrace{\begin{bmatrix}
\left[0\right]_{m \times l} \\[3pt]
\left[\mathrm{I}\right]_{l \times l} \\[3pt]
\left[0\right]_{(N-m-l)\times l}
\end{bmatrix}_{N \times l}}_{\left[\mathrm{U}\right]}
\underline{\hat{\mathrm{x}}}_{l \times 1}
+
\underbrace{\begin{bmatrix}
\left[\mathrm{I}\right]_{m \times m} & \left[0\right]_{m \times (N-m-l)} \\[3pt]
\left[0\right]_{l \times m} & \left[0\right]_{l \times (N-m-l)} \\[3pt]
\left[0\right]_{(N-m-l)\times m} & \left[\mathrm{I}\right]_{(N-m-l)\times (N-m-l)}
\end{bmatrix}_{N \times (N-l)}}_{\left[\mathrm{K}\right]}
\underline{\bar{\mathrm{x}}}_{(N-l)\times 1}
\]

The compact form becomes:

\[
\underline{\mathrm{x}} = \left[\mathrm{U}\right] \ \underline{\hat{\mathrm{x}}} + \left[\mathrm{K}\right] \ \underline{\bar{\mathrm{x}}}
\]

We remind that \(\underline{\hat{\mathrm{x}}}\) represents the unknown components, while \(\underline{\bar{\mathrm{x}}}\) represents the known components.
Notice that the unknown samples can be at arbitrary positions in the block: they do not need to form a contiguous missing segment.

We may notice that:

\[
\underline{\hat{\mathrm{x}}} = \left[\mathrm{G_x}\right]\underline{\mathrm{a}} + \underline{\mathrm{e}}
\quad \Rightarrow \quad
\underline{\mathrm{e}} = \underline{\hat{\mathrm{x}}} - \left[\mathrm{G_x}\right]\underline{\mathrm{a}}
\]

Which is explicitly:

\[
\begin{bmatrix}
e(P+1) \\[3pt]
\vdots \\[3pt]
e(N)
\end{bmatrix}
=
\begin{bmatrix}
x(P+1) \\[3pt]
\vdots \\[3pt]
x(N)
\end{bmatrix}
-
\begin{bmatrix}
x(P) & \cdots & x(1) \\[3pt]
\vdots & \ddots & \vdots \\[3pt]
x(N-1) & \cdots & x(N-P)
\end{bmatrix}
\begin{bmatrix}
a_1 \\[3pt]
\vdots \\[3pt]
a_P
\end{bmatrix}
\]

\clearpage

We can find a linear relation in $\underline{\mathrm{x}}$:

\[
\underbrace{
\begin{bmatrix}
e(P+1) \\[3pt]
\vdots \\[3pt]
e(N)
\end{bmatrix}}_{\underline{\mathrm{e}}}
=
\underbrace{
\begin{bmatrix}
-a_P & -a_{P-1} & \cdots & -a_1 & 1 & 0 & \cdots & 0 \\[3pt]
0 & -a_P & \cdots & -a_2 & -a_1 & 1 & \cdots & 0 \\[3pt]
\vdots & \vdots & \ddots & \vdots & \vdots & \vdots & \ddots & \vdots \\[3pt]
0 & 0 & \cdots & -a_P & \cdots & -a_2 & -a_1 & 1
\end{bmatrix}}_{\left[\mathrm{A}\right]}
\underbrace{
\begin{bmatrix}
x(1) \\[3pt]
x(2) \\[3pt]
\vdots \\[3pt]
x(N)
\end{bmatrix}}_{\underline{\mathrm{x}}}
\]

The conditional likelihood for white Gaussian excitation becomes:

\[
P\!\left(\underline{\hat{\mathrm{x}}} \mid \underline{\bar{\mathrm{x}}}, \ \underline{\mathrm{a}}\right)
= \left(\frac{1}{2\pi\sigma_e^2}\right)^{\tfrac{N-P}{2}}
\exp \left(-\frac{1}{2\sigma_e^2} \ \underline{\mathrm{x}}^T \left[\mathrm{A}\right]^T \ \left[\mathrm{A}\right] \ \underline{\mathrm{x}}\right)
\]

In order to express the unknown samples as a function of the known ones, we need to define an \textbf{objective function to optimize} (i.e., a cost function to minimize).
We do so by defining an estimation error, then by minimizing its energy (i.e. Least Squares method).
The cost function that we choose is the \textbf{prediction error energy} $E$ over the whole data block:

\[
E = \underline{\mathrm{e}}^{T}\,\underline{\mathrm{e}},
\qquad
\underline{\mathrm{e}} = \left[\mathrm{A}\right] \underline{\mathrm{x}}
\]

We remind that the block of data samples $\underline{\mathrm{x}}$ is assumed to be drawn from an AR process with coefficients (parameters) $a_k$.
We already showed that the excitation vector $\underline{\mathrm{e}}$ can be written as a linear function of the whole block of data samples $\underline{\mathrm{x}}$.
So, we can write:

\[
\underline{\mathrm{e}}
= \left[\mathrm{A}\right]\underline{\mathrm{x}}
= \left[\mathrm{A}\right]\left( \left[\mathrm{U}\right]\underline{\hat{\mathrm{x}}} + \left[\mathrm{K}\right]\underline{\bar{\mathrm{x}}} \right)
= \underbrace{\left[\mathrm{A}\right]\left[\mathrm{U}\right]}_{\left[\hat{A}\right]} \ \underline{\hat{\mathrm{x}}}
+ \underbrace{\left[\mathrm{A}\right]\left[\mathrm{K}\right]}_{\left[\bar{A}\right]} \ \underline{\bar{\mathrm{x}}}
= \left[\hat{A}\right]\underline{\hat{\mathrm{x}}} + \left[\bar{A}\right]\underline{\bar{\mathrm{x}}}
\]

We obtain a \textbf{column-wise partition} of $\left[\mathrm{A}\right]$ corresponding to unknown and known samples, respectively.

\begin{figure}[H]
    \centering
    \includegraphics[width=0.25\linewidth]{LSAR energy minimization.png}
    \caption{LSAR interpolation: minimization of energy $E$ with single global minimum}
\end{figure}

The LS interpolator works by \textbf{minimizing the energy} $E$ of the prediction error:

\[
\underline{\hat{\mathrm{x}}}_{\mathrm{LS}}
= \arg \min_{\underline{\hat{\mathrm{x}}}} \left\{ E \right\}
= \arg \min_{\underline{\hat{\mathrm{x}}}} \left\{ \underline{\mathrm{e}}^T \underline{\mathrm{e}} \right\}
= \arg \min_{\underline{\hat{\mathrm{x}}}} \left\{ \underline{\mathrm{x}}^T \left[\mathrm{A}\right]^T \left[\mathrm{A}\right] \, \underline{\mathrm{x}} \right\}
\]

As the prediction error is a linear function of the data, the energy $E$ is a paraboloidal function of the data.
Therefore, it has a \textbf{single global minimum}, which can be found by setting its gradient to zero:

\[
\underline{\hat{\mathrm{x}}}_{\mathrm{LS}}
= \left\{ \, \underline{\hat{\mathrm{x}}}: \ \frac{\partial E}{\partial \underline{\hat{\mathrm{x}}}} = 0 \right\}
\]

\clearpage

The computation is the following:

\[
\frac{\partial E}{\partial \underline{\hat{\mathrm{x}}}}
= 2 \underline{\mathrm{e}}^T \frac{\partial \underline{\mathrm{e}}}{\partial \underline{\hat{\mathrm{x}}}}
= 2 \underline{\mathrm{e}}^T \left[\hat{\mathrm{A}}\right]
= 2 \left( \left[\hat{\mathrm{A}}\right] \underline{\hat{\mathrm{x}}}_{\mathrm{LS}} + \left[\bar{\mathrm{A}}\right] \underline{\bar{x}} \right)^T \left[\hat{\mathrm{A}}\right]
= 0
\quad \Rightarrow \quad
\left( \underline{\hat{\mathrm{x}}}_{\mathrm{LS}}^T \left[\hat{\mathrm{A}}\right]^T + \underline{\bar{x}}^T \left[\bar{\mathrm{A}}\right]^T \right) \left[\hat{\mathrm{A}}\right]
= 0
\]

\[
\Rightarrow \quad
\underline{\hat{\mathrm{x}}}_{\mathrm{LS}}^T \left[\hat{\mathrm{A}}\right]^T \left[\hat{\mathrm{A}}\right]
= - \underline{\bar{x}}^T \left[\bar{\mathrm{A}}\right]^T \left[\hat{\mathrm{A}}\right]
\quad \Rightarrow \quad
\left[\hat{\mathrm{A}}\right]^T \left[\hat{\mathrm{A}}\right] \underline{\hat{\mathrm{x}}}_{\mathrm{LS}}
= - \left[\hat{\mathrm{A}}\right]^T \left[\bar{\mathrm{A}}\right] \underline{\bar{x}}
\]
\[
\Rightarrow \quad
\underline{\hat{\mathrm{x}}}_{\mathrm{LS}}
= - \left( \left[\hat{\mathrm{A}}\right]^T \left[\hat{\mathrm{A}}\right] \right)^{-1} \left[\hat{\mathrm{A}}\right]^T \left[\bar{\mathrm{A}}\right] \underline{\bar{x}}
\]

The method involves the solution of a set of $l$ linear equations, where $l$ is the number of missing samples in the block.  
When there are at least $P$ known samples on either side of a single contiguous burst of missing samples, the LSAR estimate can be efficiently implemented by means of Levinson–Durbin recursion (i.e. Toeplitz system of equations). 
The computational load is $\mathcal{O}(l^2)$ multiplies and additions.

\begin{figure}[H]
    \centering
    \includegraphics[width=0.4\linewidth]{Performance 1.png}
    \includegraphics[width=0.4\linewidth]{Performance 2.png}
    \caption{AR-based interpolation with $P=60$ of classical chamber music: short gaps of 25-50 samples (left), long gaps of 70-120 samples (right)}
\end{figure}

The method works well in most cases.
It performs less well for signals that do not fit the modelling assumptions, such as periodic pulse-driven voiced speech.
With longer gaps, the results tend to be dull.
However, the situation improves by increasing the order, although the loss of noisy components can be perceived.  

\clearpage
